\section*{Hoofdstuk \ref{ch:g7:sediments-and-fossils} --- Sedimenten en fossielen}

\begin{solution}{exo:g7:sediments-and-fossils:1}
(a) \omterm{def:g7:sediments-and-fossils:fossil}{afdruk}; (b) \omterm{def:g7:sediments-and-fossils:fossil}{sporenfossiel}; (c) \omterm{def:g7:sediments-and-fossils:sediment}{sediment}; (d) \omterm{def:g7:sediments-and-fossils:fossil}{afgietsel}; (e) \omterm{def:g7:sediments-and-fossils:layer}{stratum};
(f) \omterm{def:g7:sediments-and-fossils:fossil}{fossiel}.
\end{solution}

\begin{solution}{exo:g7:sediments-and-fossils:2}
(b) het dier sterft; (d) modder bedekt de schelp snel; (c) nieuwe lagen
stapelen zich op en de schelp verandert in steen; (a) \omterm{def:g5:water-shapes-the-land:erosion}{erosie} slijt de
lagen erboven weg.
\end{solution}

\begin{solution}{exo:g7:sediments-and-fossils:3}
(a) Een warme, ondiepe, heldere \omterm{def:g2:water-on-earth:ocean}{zee}. (b) Ondiep water, aan de rand van een
\omterm{def:g1:landscapes-around-us:water}{meer}, een \omterm{def:g1:landscapes-around-us:water}{rivier} of een \omterm{def:g2:water-on-earth:ocean}{zee}. (c) Een moerasbos of een moeras. (d) Land,
met groeiende planten (en vliegende libellen), bij water waar de modder
bezonk.
\end{solution}

\begin{solution}{exo:g7:sediments-and-fossils:4}
(a) Niet waar: de meeste levende wezens laten geen \omterm{def:g7:sediments-and-fossils:fossil}{fossiel} achter; ze
rotten weg of lossen op voordat ze worden begraven. (b) Waar: het is een
\omterm{def:g7:sediments-and-fossils:fossil}{sporenfossiel}. (c) Niet waar: de grootste en zwaarste korrels bezinken
eerst, de fijne modder het laatst. (d) Waar.
\end{solution}

\begin{solution}{exo:g7:sediments-and-fossils:5}
(a) Zandsteen \qty{120}{m}, kalksteen \qty{180}{m}, schalie \qty{100}{m}.
(b) $120 + 180 + 100 = \qty{400}{m}$. (c) $180 \div 400 = 0.45$:
\textbf{45~\%}. (d) De bovenkant is door \omterm{def:g5:water-shapes-the-land:erosion}{erosie} weggesleten: wat er over
is, is niet de hele laag, en hoeveel er over is, verschilt van plek tot
plek.
\end{solution}

\begin{solution}{exo:g7:sediments-and-fossils:6}
(a) Van onder naar boven: \omterm{def:g3:stones-pebbles-sand:pebble}{grind}, dan \omterm{def:g3:stones-pebbles-sand:sand}{zand}, dan modder bovenop, onder het
heldere water. (b) De piepkleine modderkorrels zinken heel langzaam: na
een uur vallen ze nog steeds. (c) Een tweede reeks lagen boven op de
eerste: weer \omterm{def:g3:stones-pebbles-sand:sand}{zand}, dan modder. (d) Het \omterm{def:g3:stones-pebbles-sand:pebble}{grind} helemaal onderaan.
\end{solution}

\begin{solution}{exo:g7:sediments-and-fossils:7}
Het best: (b), een harde schelp, snel begraven onder fijne modder in
rustig water. Dan (d): de schelp is hard en begraven, maar een \omterm{def:g5:water-shapes-the-land:flood}{overstroming}
kan haar ook verplaatsen en breken. Het minst: (a), een zacht lichaam
zonder harde delen, op rotsen waar geen modder bezinkt om het te bedekken.
(c) is niet veel beter: het bot blijft onbedekt en wordt afgebroken en
weggesleten.
\end{solution}

\begin{solution}{exo:g7:sediments-and-fossils:8}
(a) Een holte met de vorm van de schelp in de wanden gedrukt: een \omterm{def:g7:sediments-and-fossils:fossil}{afdruk}.
(b) Een \omterm{def:g7:sediments-and-fossils:fossil}{afgietsel}, een stenen kopie van de schelp. (c) Nee: de schelp is
opgelost; het \omterm{def:g7:sediments-and-fossils:fossil}{afgietsel} bestaat uit nieuwe \omterm{def:g6:rocks-and-minerals:mineral}{mineralen}. (d) Water gaat
gemakkelijk tussen de grote zandkorrels door, dus kan het de schelp
bereiken en oplossen; in fijne modder gaat het veel minder gemakkelijk.
\end{solution}

\begin{solution}{exo:g7:sediments-and-fossils:9}
(a) $1.6 \times 1000 = \qty{1600}{m}$. (b) $1600 \div 80 = 20$ lagen.
(c) $160 \div 1600 = 0.1$: \textbf{10~\%}.
\end{solution}

\begin{solution}{exo:g7:sediments-and-fossils:10}
(a) Nee: een voet laat alleen een \omterm{def:g7:sediments-and-fossils:fossil}{afdruk} achter in zachte grond, nat \omterm{def:g3:stones-pebbles-sand:sand}{zand}
of modder; de grond werd later \omterm{def:g3:stones-pebbles-sand:rock}{gesteente}. (b) De volgende golf of
\omterm{def:g5:water-shapes-the-land:flood}{overstroming} spoelt ze weg. (c) Nieuw \omterm{def:g7:sediments-and-fossils:sediment}{sediment} moest ze snel bedekken en
opvullen, voordat ze waren uitgewist. (d) Een \omterm{def:g7:sediments-and-fossils:fossil}{sporenfossiel}: het laat zien
wat het dier deed, geen deel van zijn lichaam.
\end{solution}

\begin{solution}{exo:g7:sediments-and-fossils:11}
(a) $20 + 5 = \qty{25}{mm}$. (b) $1000 \div 25 = 40$ paren: de boorkern
legt \textbf{40 \omterm{def:g5:water-shapes-the-land:flood}{overstromingen}} vast. (c) Het \omterm{def:g3:stones-pebbles-sand:sand}{zand} is zwaar: het bezinkt
kort nadat de \omterm{def:g1:landscapes-around-us:water}{rivier} het \omterm{def:g1:landscapes-around-us:water}{meer} binnenkomt, bij de \omterm{def:g5:water-shapes-the-land:source-mouth}{monding}; maar een klein
deel ervan komt ver. (d) Ver weg zit alleen de fijne modder nog in het
water: die bezinkt daar langzaam, als het water rustig is.
\end{solution}

\begin{solution}{exo:g7:sediments-and-fossils:12}
(a) Schelpen zijn hard: ze blijven lang genoeg bewaard om te worden
begraven, en veel dieren met schelpen leefden op de zeebodem, waar modder
over hen heen bezinkt. (b) Zachte lichamen rotten weg voordat ze kunnen
worden begraven. (c) Nee: zachte dieren laten bijna nooit \omterm{def:g7:sediments-and-fossils:fossil}{fossielen}
achter, dus de lijst uit de \omterm{def:g7:sediments-and-fossils:fossil}{fossielen} is onvolledig; ze toont vooral de
dieren met schelpen of botten.
\end{solution}

\begin{solution}{pb:g7:sediments-and-fossils:1}
\textbf{1.} Zandsteen, kalksteen, schalie, krijt (de eerste onderaan de
kolom).

\textbf{2.} \qty{120}{m}, \qty{180}{m} en \qty{100}{m}; samen
\qty{400}{m}.

\textbf{3.} $180 \div 400 = 0.45$: \textbf{45~\%}.

\textbf{4.} De bovenkant is weggesleten: het bovenste deel van de laag is
verdwenen, en niemand kan meten wat er niet meer is.

\textbf{5.} In ondiep water aan de rand van een \omterm{def:g1:landscapes-around-us:water}{meer}, een \omterm{def:g1:landscapes-around-us:water}{rivier} of een
\omterm{def:g2:water-on-earth:ocean}{zee}: golfribbels ontstaan in \omterm{def:g3:stones-pebbles-sand:sand}{zand} in ondiep water, en voetafdrukken worden
achtergelaten in zacht, nat \omterm{def:g3:stones-pebbles-sand:sand}{zand} of modder, vaak aan een \omterm{def:g1:landscapes-around-us:water}{kust}.

\textbf{6.} In een warme, ondiepe, heldere \omterm{def:g2:water-on-earth:ocean}{zee}: de koralen.

\textbf{7.} Rustig: fijne modder bezinkt alleen waar het water bijna stilstaat.

\textbf{8.} Eerst was de plek een \omterm{def:g1:landscapes-around-us:water}{kust}, met ondiep water en zacht \omterm{def:g3:stones-pebbles-sand:sand}{zand}
waarover dieren liepen. Daarna bedekte de \omterm{def:g2:water-on-earth:ocean}{zee} haar: een warme, ondiepe \omterm{def:g2:water-on-earth:ocean}{zee}
waarin koralen groeiden en schelpen zich opstapelden. Later werd het water
rustiger, misschien dieper, en bezonk er alleen fijne modder, met de
schelpen van ammonieten, dieren die verwant zijn aan inktvissen en
octopussen.

\textbf{9.} $180 \times 100 = 18\,000$~cm.

\textbf{10.} $120~\text{m} = 12\,000$~cm, en $12\,000 \times 20 =
240\,000$~jaar.

\textbf{11.} $1~\text{m} = 100$~cm, en $100 \times 100 = 10\,000$~jaar.

\textbf{12.} $180 \times 10\,000 = 1\,800\,000$: ongeveer
\textbf{1\,800\,000 jaar}, bijna twee miljoen jaar. Dat is
$1\,800\,000 \div 240\,000 = 7.5$ keer langer dan voor het zandsteen.
\end{solution}
